% УВОД, без номер. Ориентировъчен обем: 1 до 2 стр.
% Резюме на дейностите по разработката на курсовия проект.
\section*{УВОД}
\label{sec:introduction}
\addcontentsline{toc}{section}{УВОД}

Маршрутизацията в мрежите по стека TCP/IP се описва в учебната литература като решена
задача: всеки маршрутизатор поддържа представа за топологията, изчислява най-кратките
пътища и попълва маршрутна таблица. Поведението, което определя качеството на мрежата
на практика, обаче не е в установеното състояние, а в прехода между две установени
състояния. Когато една връзка откаже, различните възли научават за отказа в различен
момент и за краткия интервал, докато представите им се разминават, мрежата може да
образува преходни цикли и да губи пакети по пътища, които изглеждат валидни за всеки
възел поотделно. Този интервал се нарича време за сходимост и е основният предмет на
настоящия курсов проект.

Причината преходът да се изучава рядко в лабораторни условия е практическа. Отказът на
връзка в реална мрежа настъпва в момент, който не се повтаря, а таймерите на протокола
се задвижват от истински часовник. Две последователни наблюдения на едно и също събитие
дават различни числа и разликата между тях не може да бъде отдадена на промяна в
протокола. За да бъде измерим, преходът трябва да бъде възпроизводим, а
възпроизводимостта изисква времето да бъде управлявано от програмата, а не от часовника
на машината.

\textbf{Целта} на проекта е да се изгради демон за маршрутизация по състояние на
връзките от семейството OSPF заедно с дискретно-събитиен симулатор, който изпълнява
същия програмен код върху управлявано време, и с този инструмент да се измери времето
за сходимост и образуването на преходни цикли при отказ на връзка.

За постигането на целта са формулирани следните \textbf{задачи}:

\begin{enumerate}[topsep=0pt,itemsep=0pt,leftmargin=*]
\item литературен преглед на принципите на маршрутизация и управление на потоците в
      мрежи по TCP/IP, поставен в контекста на дисциплината;
\item анализ на възможните решения по езика, модела на времето, транспортния слой и
      еталонната реализация, и обосновка на направения избор;
\item проектиране на демона така, че протоколната логика да не зависи от източника на
      време и от източника на пакети;
\item програмна реализация на откриване на съседи, разпространение на обявленията за
      състояние на връзките, изчисляване на най-кратките пътища и инсталиране на
      маршрути;
\item изграждане на методика за възпроизводимо измерване на сходимостта при отказ на
      връзка;
\item провеждане на измерванията и анализ на получените резултати.
\end{enumerate}

\textbf{Обхватът} на реализацията е маршрутизиращият слой на стека TCP/IP. Безжичните,
спътниковите и клетъчните мрежи, както и ATM, присъстват в раздел~\ref{sec:literature}
като контекст, защото са част от учебното съдържание на дисциплината, но не влизат в
реализацията. Причината е методологична, а не липса на интерес: тези среди въвеждат
модел на канала със загуби и променлив капацитет, чиято достоверност сама по себе си е
отделна изследователска задача, и включването им би направило измерената сходимост
зависима от достоверността на модела на канала вместо от поведението на протокола.

\textbf{Структура на изложението.} Раздел~\ref{sec:literature} представя предметната
област. Раздел~\ref{sec:alternatives} излага разгледаните алтернативи и обосновава
избора. Раздел~\ref{sec:design} описва проектирането, а раздел~\ref{sec:implementation}
програмната реализация. Раздел~\ref{sec:method} описва методиката за измерване, а
раздел~\ref{sec:results} получените резултати. Заключението обобщава изводите и
очертава посоките за по-нататъшна работа.
